\documentclass[11pt]{article}
\usepackage[T1]{fontenc}
\usepackage[utf8]{inputenc}
\usepackage{lmodern}
\usepackage{amsmath,amssymb}
\usepackage{longtable,booktabs,array}
% Compatibility with Pandoc's captionless tables on older longtable versions.
\ifcsname c@none\endcsname\else\newcounter{none}\fi
\usepackage{graphicx}
\usepackage[margin=25mm]{geometry}
\usepackage{hyperref}
\hypersetup{colorlinks=true,linkcolor=blue,urlcolor=blue}
\providecommand{\tightlist}{\setlength{\itemsep}{0pt}\setlength{\parskip}{0pt}}
\providecommand{\pandocbounded}[1]{#1}
\setlength{\emergencystretch}{3em}
\setcounter{secnumdepth}{-1}
\begin{document}
\section{Grothendieck's existence
theorem}\label{grothendiecks-existence-theorem}

\emph{Written by GPT-6.1 Sol (OpenAI), in Codex, at Ultra effort,
October 2026. Self-checked by the writing AI, GPT-6.1 Sol, at Ultra
effort. Public domain (CC0).}

Formal functions recovers completed cohomology from infinitesimal
neighborhoods. Grothendieck's existence theorem recovers the sheaves
themselves. Over a complete Noetherian base, a compatible coherent sheaf
on every thickening of a proper scheme comes from one coherent sheaf,
and compatible morphisms come from unique algebraic morphisms.

We use \href{the-theorem-on-formal-functions.md}{formal functions},
\href{serres-theorems-on-projective-schemes.md}{Serre generation and
vanishing}, and
\href{proper-morphisms-and-coherent-direct-images.md}{proper coherence}.
The planned \emph{Completion} lesson in \emph{Commutative algebra for
geometry} supplies Artin--Rees, Noetherianness and flatness of
completion, and completeness of finite modules over a complete
Noetherian ring. Chow's lemma belongs to the planned \emph{Projective
morphisms and Chow's lemma} lesson in \emph{Morphisms of schemes}; its
exact open proof was specified in our proper-coherence lesson.

\subsection{1. Compatible coherent
systems}\label{compatible-coherent-systems}

For a Noetherian scheme \(X\) and a coherent ideal
\(J\subset\mathcal O_X\), a \textbf{coherent formal module} is a system
\(\mathfrak F=(F_n)_{n\geq1}\) of coherent sheaves, with \[
J^nF_n=0,\qquad F_{n+1}/J^nF_{n+1}\xrightarrow{\sim}F_n.
\tag{1}
\] Morphisms are compatible sheaf maps at every level. Equivalently,
\(F_n\) is a coherent sheaf on \(X_n=V(J^n)\). Iterating (1) gives
\(F_m/J^nF_m=F_n\) for \(m\geq n\). Write \(\operatorname{Coh}(X,J)\)
for this category.

\textbf{Lemma 1.1 (the affine description).} If
\(X=\operatorname{Spec}R\) and \(J=I\mathcal O_X\), this category is
equivalent to finite modules over \(\widehat R=\varprojlim R/I^n\). The
correspondence is \[
(M_n)\longmapsto M=\varprojlim M_n,
\qquad M\longmapsto(M/I^nM).
\]

\textbf{Proof.} The transition maps are onto, so a compatible sequence
lifting any chosen element of \(M_1\) can be constructed recursively.
Lift finitely many generators of \(M_1\) to \(M\). They define maps
\((R/I^n)^r\to M_n\), which are onto by Nakayama for the nilpotent ideal
\(I/I^n\).

Let \(K_n\) be their kernels. These kernels also have surjective
transitions. To lift \(v\in K_n\), lift it to \(w\in(R/I^{n+1})^r\). Its
image in \(M_{n+1}\) lies in \(I^nM_{n+1}\), since it vanishes in
\(M_n\). Surjectivity of the free-module map lets us subtract an element
of \(I^n(R/I^{n+1})^r\) with that same image. The corrected lift belongs
to \(K_{n+1}\).

Taking limits gives \[
0\to\varprojlim K_n\to\widehat R^{\,r}\to M\to0.
\] Surjectivity on the right follows by adjusting consecutive lifts
using the surjective kernel transitions. Thus \(M\) is finite. Since
\(\varprojlim K_n\to K_n\) is onto, reducing this presentation modulo
\(I^n\) gives \(M/I^nM=M_n\). Conversely finite modules over the
Noetherian complete ring \(\widehat R\) are complete, so are recovered
from their quotients. Compatible maps pass uniquely to and from limits,
proving the equivalence. \(\square\)

An affine open in a scheme over a complete ring need not have a complete
coordinate ring. For instance the \(t\)-adic completion of \(k[[t]][x]\)
is the ring \[
k[[t]]\langle x\rangle=\varprojlim_n(k[t]/t^n)[x]
\] of restricted power series: in \(\sum a_jx^j\), the coefficients
\(a_j\) tend to zero \(t\)-adically. This contains series of unbounded
degree in \(x\), unlike the polynomial ring.

\subsection{2. Exactness and the actual formal
kernel}\label{exactness-and-the-actual-formal-kernel}

\textbf{Theorem 2.1.} The category \(\operatorname{Coh}(X,J)\) is
abelian, exactness is local on \(X\), and the completion functor \[
\operatorname{Coh}(X)\longrightarrow\operatorname{Coh}(X,J),
\qquad F\longmapsto\widehat F=(F/J^nF)
\] is exact.

\textbf{Proof.} On an affine, Lemma 1.1 identifies the category with
finite modules over a Noetherian ring, an abelian category. We explain
why its kernels glue, since taking the kernels separately at each level
is generally wrong.

For a morphism \(\alpha:(F_n)\to(G_n)\), take the images \[
K'_{l,m}=\operatorname{im}(\ker\alpha_l\to F_m),\qquad l\geq m.
\] On an affine let \(u:M\to N\) be the corresponding completed-module
map, with kernel \(K\) and image \(H\). Artin--Rees for \(H\subset N\)
gives a constant \(c\) such that \[
u^{-1}(I^lN)\subset K+I^{l-c}M\quad(l\geq c).
\] Indeed \(H\cap I^lN\subset I^{l-c}H\), and elements of this last
module lift from \(I^{l-c}M\). Hence for \(l\geq m+c\), the image
\(K'_{l,m}\) is precisely \(K/(K\cap I^mM)\). It stabilizes. Artin--Rees
for \(K\subset M\) then shows that, for fixed \(n\) and sufficiently
large \(m\), \[
K'_m/I^nK'_m
=K/(K\cap I^mM+I^nK)=K/I^nK.
\tag{2}
\] All these images and quotients are constructions with coherent
sheaves and commute with restriction. A finite affine cover gives
uniform stabilization bounds, so they define coherent sheaves globally.
The resulting system is the formal kernel: a map of systems killed by
\(\alpha\) factors at every level through the stabilized images, then
through (2); the factorization is unique. Cokernels are the levelwise
cokernels, since quotient formation preserves (1). On affines these
constructions give the ordinary module kernel and cokernel, so image
equals coimage. They also show that restriction is exact and that
exactness can be checked on an open cover.

On an affine the completion functor is \(M\mapsto M\otimes_R\widehat R\)
on finite modules. Flatness of Noetherian completion makes it exact. The
local characterization just established proves exactness globally.
\(\square\)

For example multiplication by \(t\) on \(k[[t]]\) has zero kernel, but
multiplication by \(t\) on \(k[t]/t^n\) has kernel \(k\,t^{n-1}\). These
nonzero kernels map to zero at the preceding level. Their stabilized
formal kernel is zero. An exact sequence in the formal category
therefore need not give a left-exact sequence at each fixed quotient
level.

A map of formal modules is onto exactly when its first-level map is
onto. In the affine description, its cokernel is a finite completed
module \(C\); if \(C/IC=0\), Nakayama gives \(C=0\). Here
\(I\widehat R\) lies in the Jacobson radical: \(1-a\) is invertible for
\(a\in I\widehat R\), by its convergent geometric series. This also
proves the assertion locally on \(X\).

\subsection{3. Full faithfulness from formal
functions}\label{full-faithfulness-from-formal-functions}

Now let \(A\) be Noetherian and complete for an ideal \(I\), let \(X\)
be proper over \(A\), and set \(J=I\mathcal O_X\).

\textbf{Lemma 3.1 (completed Hom).} For coherent \(F,G\), with
\(H=\mathcal Hom(F,G)\), \[
\operatorname{Hom}_{\operatorname{Coh}(X,J)}(\widehat F,\widehat G)
\cong\varprojlim_n\Gamma(X,H/J^nH).
\tag{3}
\]

\textbf{Proof.} Affine locally, \(F\) has a finite presentation.
Applying \(\operatorname{Hom}(-,G)\) describes Hom as the kernel of a
map between two finite direct sums of \(G\). Tensoring this kernel with
the flat completed ring preserves the kernel and gives \[
\operatorname{Hom}_R(M,N)\otimes_R\widehat R
=\operatorname{Hom}_{\widehat R}(\widehat M,\widehat N).
\] Lemma 1.1 identifies the right side with compatible maps on the
quotients. These natural affine identifications identify the sheaves of
completed Hom sections with the sheaves of compatible formal maps. They
glue on overlaps. Taking global sections, which commutes with inverse
limits of sheaves, gives (3). Notice that no equality of the separate
finite-level Hom modules was asserted. \(\square\)

\textbf{Theorem 3.2 (full faithfulness).} Completion on coherent sheaves
on \(X\) is fully faithful.

\textbf{Proof.} The sheaf \(H\) is coherent, so its global section
module is finite over \(A\). Formal functions identifies the right side
of (3) with \(\widehat{\Gamma(X,H)}\). A finite module over the complete
Noetherian ring \(A\) is already complete. Thus (3) identifies formal
maps with \(\Gamma(X,H)=\operatorname{Hom}_X(F,G)\), and the
identification is the completion map on morphisms. \(\square\)

In particular, an algebraization, when it exists, is unique up to the
unique isomorphism inducing a specified formal identification. Formal
isomorphisms algebraize because both directions algebraize and
faithfulness makes their composites identities.

\subsection{4. One twist for every
thickening}\label{one-twist-for-every-thickening}

Suppose \(X\) is projective over \(A\), and choose a relatively ample
invertible sheaf \(L\). For a formal module \(\mathfrak F=(F_n)\), put
\(F_0=0\) and \[
D_n=\ker(F_{n+1}\to F_n),\qquad n\geq0.
\]

\textbf{Lemma 4.1 (uniform vanishing).} There exists \(d_0\) such that
\[
H^1(X,D_n\otimes L^d)=0\quad\text{for every }n\geq0,
\quad d\geq d_0.
\]

\textbf{Proof.} Let \(B=\bigoplus_{n\geq0}I^n/I^{n+1}\), a Noetherian
graded algebra generated in degree one over \(A/I\). The sheaf
\(D=\bigoplus D_n\) is a graded module over \(\mathcal O_X\otimes_A B\).
On an affine, Lemma 1.1 identifies it with the associated graded module
\(\bigoplus J^nM/J^{n+1}M\) of a finite completed module. It is
generated over the associated graded ring in degree zero. The map from
the pulled-back \(B\) to that associated graded ring is surjective,
because \(J\) is generated by the image of \(I\). Consequently \(D\) is
of finite type over \(\mathcal O_X\otimes_A B\). Surjectivity is
sufficient; no flatness is needed to identify associated graded rings.

On \(X_B=X\times_A\operatorname{Spec}B\), with affine projection \(p\),
the affine sheaf correspondence realizes \(D=p_*E\) for a coherent
\(E\). The morphism \(X_B\to\operatorname{Spec}B\) is projective and
\(p^*L\) is relatively ample. Serre vanishing gives one bound for the
positive cohomology of \(E\otimes p^*L^d\). Affine pushforward, its
projection formula, and cohomology commuting with direct sums of
quasi-coherent sheaves on this separated quasi-compact scheme identify
that cohomology with \(\bigoplus_n H^1(X,D_n\otimes L^d)\). Every
summand is zero for the same bound. \(\square\)

The bound is uniform in the infinitesimal index. Applying Serre
vanishing independently to each \(F_n\) would give bounds depending on
\(n\), insufficient for a single finite presentation.

\subsection{5. Existence on a projective
scheme}\label{existence-on-a-projective-scheme}

\textbf{Theorem 5.1 (projective existence).} For projective \(X\) over
\(A\), completion gives an equivalence \[
\operatorname{Coh}(X)\simeq\operatorname{Coh}(X,I\mathcal O_X).
\]

\textbf{Proof.} Full faithfulness is Theorem 3.2. Let \(\mathfrak F\) be
a formal module. Choose \(d\) large enough both for Lemma 4.1 and for
finitely many sections of \(F_1\otimes L^d\) to generate it. The exact
sequences with kernel \(D_n\otimes L^d\) make \[
\Gamma(X,F_{n+1}\otimes L^d)\to\Gamma(X,F_n\otimes L^d)
\] surjective. Lift the chosen finite collection recursively to
compatible sections at every level. They give a formal morphism \[
\widehat P\twoheadrightarrow\mathfrak F,
\qquad P=(L^{-d})^r.
\] It is onto because it is onto at the first level. Its kernel
\(\mathfrak K\) is a coherent formal module by Theorem 2.1. Apply the
same construction to \(\mathfrak K\) to obtain
\(\widehat Q\twoheadrightarrow\mathfrak K\), with \(Q=(L^{-e})^s\) for
some \(e,s\). Thus \[
\widehat Q\longrightarrow\widehat P\longrightarrow\mathfrak F\to0
\] is a presentation in the formal category. Full faithfulness
algebraizes its first map to \(Q\to P\). Take its coherent cokernel
\(F\). Exactness of completion identifies \(\widehat F\) with
\(\mathfrak F\), proving essential surjectivity. \(\square\)

The two presentations use different twists if necessary. There is no
assertion that their kernel is a vector bundle, or that every coherent
sheaf has a finite resolution by line bundles.

\subsection{6. Proper schemes and proper
supports}\label{proper-schemes-and-proper-supports}

\textbf{Theorem 6.1 (Grothendieck existence).} Let \(A\) be Noetherian
and \(I\)-adically complete, and let \(X\) be separated and of finite
type over \(A\). Completion gives an equivalence between coherent
sheaves on \(X\) with support proper over \(A\) and coherent formal
modules whose first-level support is proper over \(A\). In particular,
when \(X\) is proper, all coherent sheaves and all coherent formal
modules participate.

The general reduction has a precise open proof in {[}Stacks, Tags 088C
and 088E{]}. We explain how it extends Theorem 5.1, including the two
technical descent steps and their proof providers. This uses the theorem
with proper supports in its full stated generality.

For proper \(X\), argue by Noetherian induction on closed subschemes: a
maximal ideal defining a counterexample permits assuming existence for
every closed subscheme defined by a nonzero ideal. Chow's lemma gives a
proper surjection \(p:Y\to X\), an isomorphism over a dense open \(U\),
with \(Y\) projective over \(A\). Pull the formal module back to \(Y\)
and algebraize it by Theorem 5.1, obtaining \(E\). Proper coherence
makes \(H=p_*E\) coherent.

The precise comparison lemma {[}Stacks, Tag 088B{]} constructs \[
\alpha:\mathfrak F\longrightarrow\widehat H
\] whose kernel and cokernel are killed by a fixed power of the ideal
\(K\) of \(X\setminus U\). Its construction takes the inverse limit of
the adjunction maps to \(p_*(E/I^nE)\); formal functions identifies this
limit with the completed pushforward. To establish the fixed bound, work
over an affine of \(X\), change flatly to its complete coordinate ring,
and apply full faithfulness for the proper morphism there. The
comparison becomes the completion of an actual adjunction map, which is
an isomorphism over \(U\). Coherent kernel and cokernel supported on
\(V(K)\) are killed by one power of \(K\).

The second technical step {[}Stacks, Tag 088A{]} says that such a
bounded comparison algebraizes \(\mathfrak F\), provided existence holds
on every \(V(K^e)\). It algebraizes \(\mathfrak F/K^e\mathfrak F\),
changes the direction of completion from \(I\) to \(K\), and uses
bounded-error formal gluing {[}Stacks, Tag 0889{]}. The uniqueness there
recovers the original \(I\)-quotients. The induction hypothesis supplies
existence on all \(V(K^e)\): the dense open meets every component, so
\(K^e\neq0\). These steps, with their exact proofs, complete the proper
case in {[}Stacks, Tag 088C{]}. This is an ideal induction, requiring no
finite bound on the dimension of the Noetherian base.

For the proper-support assertion, the modification \(Y\) is only
quasi-projective over \(A\). Embed it openly in a projective
\(\overline Y\). The support of each pulled-back \(F_n\) equals the
inverse image of the support of \(F_1\), hence is proper. Its image in
\(\overline Y\) is closed, so extension by zero is coherent and gives a
formal module on \(\overline Y\). Algebraize this by Theorem 5.1. The
algebraized support cannot meet the boundary: its boundary intersection
is closed in the projective scheme, has closed image in
\(\operatorname{Spec}A\), and misses \(V(I)\). But \(I\) lies in the
Jacobson radical of the complete ring \(A\), so every nonempty closed
subset of \(\operatorname{Spec}A\) meets \(V(I)\). Thus the support lies
in \(Y\) and is proper.

Push forward and use the same comparison. Noetherian induction and the
Serre-category property of formal modules with proper support show that
the original formal module is annihilated by the ideal of some closed
proper subscheme \(Z\subset X\). On \(Z\), the proper case applies.
Likewise every coherent sheaf with proper support lives on its proper
scheme-theoretic support. Hom maps between any two such sheaves can be
computed on the proper closed union of their supports, so full
faithfulness follows there. This is the support reduction and full proof
furnished by {[}Stacks, Tag 088E{]}. Properness of \(X\) itself has not
been inserted into that statement.

\subsection{7. Examples: what can fail on an affine
line}\label{examples-what-can-fail-on-an-affine-line}

On \(\mathbf P^1_{k[[t]]}\), the systems \(\mathcal O_{X_n}(d)\)
algebraize to \(\mathcal O_X(d)\). The extension system \[
0\to\mathcal O_{X_n}(-2)\to E_n\to\mathcal O_{X_n}\to0
\] with compatible class \(a_n\in k[t]/t^n\) algebraizes to the
extension with class \(a=\varprojlim a_n\in k[[t]]\), because
\(H^1(\mathcal O_X(-2))=k[[t]]\). Choosing \(a=t\) recovers our
cohomology-jump example.

For \(X=\mathbf A^1_{k[[t]]}\), completion is not fully faithful. The
restricted series \[
h(x)=\sum_{j\geq0}t^jx^j
\] defines a compatible multiplication map on the systems
\(\mathcal O_{X_n}\), but is not a polynomial in \(x\). It therefore
does not come from an endomorphism of \(\mathcal O_X\). Faithfulness
also fails: the nonzero coherent module \(k[[t]][x]/(tx-1)\) has zero
completion, since \(tx-1\) is a unit modulo every power of \(t\).

These are failures of morphism recovery and uniqueness, not examples of
nonalgebraizable objects. In fact this particular affine line has a
stronger existence property.

\textbf{Proposition 7.1.} Every coherent formal module on
\(\mathbf A^1_{k[[t]]}\) is algebraizable.

\textbf{Proof.} Let \(R=k[[t]][x]\), \(B=\widehat R\), and let the
system correspond to a finite \(B\)-module \(M\). The ring \(B\) is a
Noetherian domain, and \(B/tB=k[x]\). Thus \((t)\) is prime, and
\(B_{(t)}\) is a discrete valuation ring: its maximal ideal is generated
by the nonzero element \(t\). The structure theorem over this ring gives
\[
M_{(t)}\cong B_{(t)}^{\,r}\oplus\bigoplus_{i=1}^s B_{(t)}/(t^{a_i}).
\] Finite presentations let this isomorphism and its inverse be defined
after inverting finitely many elements of \(B\setminus(t)\). The
identities defining inverse maps also hold after finitely many such
inversions. Let \(h\in k[x]\setminus\{0\}\) be the product of their
nonzero reductions modulo \(t\). In the \(t\)-adic completion of
\(R[1/h]\), each denominator is a unit, since its reduction modulo \(t\)
is a unit. Hence the formal system on \(D(h)\) is isomorphic to the
completion of the displayed free-plus-\(t\)-power-torsion module.

The zero set of \(h\) in \(\mathbf P^1_k\) is finite and lies in the
affine line: the homogenization has nonzero value at infinity. Let
\(V\subset\mathbf P^1_{k[[t]]}\) be the complement of that homogenized
zero set. It contains infinity and meets the affine chart in \(D(h)\).
On \(V\), take the formal completion of the same free-plus-torsion
module. Glue it to the given system on the affine chart using the formal
isomorphism on \(D(h)\). This gives compatible coherent sheaves on all
projective-line thickenings. Theorem 5.1 algebraizes them on
\(\mathbf P^1_{k[[t]]}\). Restricting that sheaf to the affine chart
algebraizes the original system. \(\square\)

Thus a claim that an arbitrary formal multiplication series
automatically gives a nonalgebraizable coherent module on this affine
line would be false. The theorem's equivalence, including morphisms and
unique recovery, genuinely needs its support hypotheses; essential
surjectivity alone can hold in a special nonproper situation.

\subsection{8. Exercises with solutions}\label{exercises-with-solutions}

\textbf{Exercise 8.1 (easy: exactness).} Explain the sense in which
completion of a coherent short exact sequence is exact, and exhibit the
failure of left exactness at a fixed quotient level.

\textbf{Solution.} On an affine, finite-module completion is tensoring
with the flat Noetherian completed ring, so is exact; Theorem 2.1 glues
this statement in the formal category. For multiplication by \(t\) on
\(k[[t]]\), the original map is injective but its reduction modulo
\(t^n\) has kernel \(k t^{n-1}\). These kernels are transient and have
stabilized formal kernel zero. Thus finite-level injectivity is not the
meaning of exact completion.

\textbf{Exercise 8.2 (medium: projective-line maps).} Compute the
compatible morphisms from \((\mathcal O_{X_n}(a))\) to
\((\mathcal O_{X_n}(b))\) on \(\mathbf P^1_{k[[t]]}\).

\textbf{Solution.} They are the limit of sections of
\(\mathcal O_{X_n}(b-a)\). If \(b-a<0\), each group is zero. Otherwise
its monomial basis gives \((k[t]/t^n)^{b-a+1}\), whose limit is
\(k[[t]]^{b-a+1}\). This is exactly the algebraic Hom group from
projective-space cohomology, with its coefficientwise restriction map.
It proves full faithfulness for these line bundles directly.

\textbf{Exercise 8.3 (medium: nonproper failure).} Give both a
compatible endomorphism on the affine-line formal structure sheaf that
does not algebraize as an endomorphism of \(\mathcal O_X\), and a
nonzero coherent sheaf with zero completion.

\textbf{Solution.} Multiplication by \(\sum_{j\geq0}t^jx^j\) is defined
modulo every \(t^n\), but a polynomial cannot have its unbounded nonzero
\(x\)-coefficients, so it does not algebraize as the specified
endomorphism. The sheaf associated to \(R/(tx-1)\) is nonzero: its ring
is \(k((t))\), with \(x=t^{-1}\). Its reduction modulo \(t^n\) is zero
because \(t\) is invertible there. Proposition 7.1 explains why neither
example should be misdescribed as a nonalgebraizable formal object.

\textbf{Exercise 8.4 (medium: algebraizing a presentation).} Recover
projective existence from uniform vanishing, full faithfulness and
exactness.

\textbf{Solution.} Use a uniform twist to lift a finite generating
collection for \(F_1\), giving
\(\widehat P\twoheadrightarrow\mathfrak F\). Its formal kernel is
coherent. Do the same for that kernel to obtain
\(\widehat Q\to\widehat P\to\mathfrak F\to0\). Full faithfulness
algebraizes \(Q\to P\), and exactness identifies the completion of its
coherent cokernel with \(\mathfrak F\). The construction never assumes
the finite-level kernels themselves form the formal kernel.

\textbf{Exercise 8.5 (hard: the proper reduction).} Identify the exact
two descent inputs needed after algebraizing a formal pullback on a Chow
modification, and check their hypotheses.

\textbf{Solution.} The proper modification is an isomorphism away from
\(V(K)\); its algebraized pullback is coherent. Tag 088B applies to
these hypotheses and gives a map to the completed coherent pushforward
with kernel and cokernel killed by a fixed power of \(K\). Tag 088A
applies once existence is known on every \(V(K^e)\), and turns this
bounded comparison into algebraization. Noetherian induction supplies
those closed-subscheme equivalences, with full faithfulness on them
supplied by Theorem 3.2. The exact open proofs linked in Section 6
furnish the comparison and the change-of-completion gluing step; this is
the complete reduction used by Tag 088C.

\textbf{Exercise 8.6 (hard: a special affine existence theorem).}
Explain the two features of the affine line that make Proposition 7.1
work.

\textbf{Solution.} At the generic point of its closed fiber, the
completed ring localized at \((t)\) is a discrete valuation ring. A
finite module there has a free-plus-\(t\)-power-torsion normal form,
which spreads to \(D(h)\) after completing, using finitely many
denominators with nonzero reductions. The zero set of the
single-variable polynomial \(h\) has projective closure disjoint from
infinity. Thus the standard module can be put on an open neighborhood of
the entire boundary, and glued to the original system. Projective
existence finishes. In higher dimension a hypersurface's closure can
meet the boundary, so this argument does not establish unrestricted
affine existence in general.

\subsection{References}\label{references}

\begin{itemize}
\tightlist
\item
  \textbf{{[}Stacks{]}} The Stacks project authors, \emph{The Stacks
  project}, in its AI Integrated Stacks Project edition: affine formal
  modules
  \href{https://kokunoyumeto.github.io/stacks-zh-hans-cn/en/coherent.html\#coherent-lemma-inverse-systems-affine}{Tag
  087W}, their abelian category
  \href{https://kokunoyumeto.github.io/stacks-zh-hans-cn/en/coherent.html\#coherent-lemma-inverse-systems-abelian}{Tag
  087X}, exact completion
  \href{https://kokunoyumeto.github.io/stacks-zh-hans-cn/en/coherent.html\#coherent-lemma-exact}{Tag
  0881}, completed Hom
  \href{https://kokunoyumeto.github.io/stacks-zh-hans-cn/en/coherent.html\#coherent-lemma-completion-internal-hom}{Tag
  0882}, full faithfulness
  \href{https://kokunoyumeto.github.io/stacks-zh-hans-cn/en/coherent.html\#coherent-lemma-fully-faithful}{Tag
  0883}.
\item
  Uniform vanishing
  \href{https://kokunoyumeto.github.io/stacks-zh-hans-cn/en/coherent.html\#coherent-lemma-vanishing-projective}{Tag
  0884}, and projective existence
  \href{https://kokunoyumeto.github.io/stacks-zh-hans-cn/en/coherent.html\#coherent-lemma-existence-projective}{Tag
  0885}.
\item
  The exact open proper-case descent proofs are the bounded comparison
  \href{https://kokunoyumeto.github.io/stacks-zh-hans-cn/en/coherent.html\#coherent-lemma-inverse-systems-push-pull}{Tag
  088B}, change of completion
  \href{https://kokunoyumeto.github.io/stacks-zh-hans-cn/en/coherent.html\#coherent-lemma-existence-tricky}{Tag
  088A}, and bounded formal gluing
  \href{https://kokunoyumeto.github.io/stacks-zh-hans-cn/en/coherent.html\#coherent-lemma-existence-easy}{Tag
  0889}. They assemble into proper existence
  \href{https://kokunoyumeto.github.io/stacks-zh-hans-cn/en/coherent.html\#coherent-proposition-existence-proper}{Tag
  088C} and proper-support existence
  \href{https://kokunoyumeto.github.io/stacks-zh-hans-cn/en/coherent.html\#coherent-theorem-grothendieck-existence}{Tag
  088E}.
\item
  These linked reference proofs retain GNU FDL 1.2. Sections 1--5, the
  affine-line argument, examples and solutions are independently written
  CC0; Section 6 uses the specified open technical proof providers in
  the full generality of the stated theorem.
\end{itemize}
\end{document}
